\section{Diagnostic Failure Modes}\label{sec:failuremodes}\label{sec:examples}

The decomposition $(\Pot,\Act,\Conv)$ becomes useful only when it separates
failure regimes that matter empirically. This section names the three failure
modes that recur throughout the experiments and later feed directly into the
coherence scores. The point is not to enumerate every possible pathology of an
adaptive system, but to isolate the three ways in which representation,
actuation, and adaptation can fail in the diagnostic story developed here.

\begin{proposition}[Failure Mode~1: Static-$\Pot$ collapse]\label{prop:failP}
Suppose $\Conv_t$ ceases to update $\Pot$, so $\Pot(S,t)=\Pot(S,t_0)$ for
all $t>t_0$. If the environment undergoes distributional shift
$\Pot^*(S,t)\neq\Pot^*(S,t_0)$ for large $t$, then
$W_2(\Pot(S,t),\Pot^*(S,t))\to\infty$ and the expected actuation error
$\mathbb{E}[d_\mathcal{O}(\Act_t(\Pot),\Act^*_t)]$ is unbounded.
\end{proposition}

\begin{proof}
$\Pot(S,t)=\Pot(S,t_0)$ is fixed while $\Pot^*$ drifts, so $W_2\to\infty$ by
definition. Lipschitz continuity of $\Act_t$ then propagates the divergence to
actuation error.
\end{proof}

\begin{remark}[Interpretation]
FM-1 is the stale-model regime: the representation no longer tracks the world,
so downstream actions inherit an increasingly obsolete state estimate.
\end{remark}

\begin{proposition}[Failure Mode~2: Non-deterministic-$\Act$ divergence]\label{prop:failA}
If $\Act_t$ is not $\delta$-reproducible for any $\delta<\infty$, then
$H(\Act_t \mid \Pot(S,t))>0$ for all $t$: actuation carries information not
determined by the state model. Distributed replicas of the system diverge, and
the composite resists auditing and replication.
\end{proposition}

\begin{proof}
Non-$\delta$-reproducibility implies $\exists\,\omega\neq\omega'$ with
$\Act_t(p;\omega)\neq\Act_t(p;\omega')$, giving positive conditional entropy.
Auditability requires reconstructing actuation from state; this fails when
conditional entropy is positive.
\end{proof}

\begin{remark}[Interpretation]
FM-2 is the non-reproducible-action regime: the model may still be accurate,
but the realized action can no longer be reconstructed from the represented
state, so replicas and audits diverge.
\end{remark}

\begin{proposition}[Failure Mode~3: Absent-$\Conv$ monopolization]\label{prop:failPhi}
Let $\Conv_t\equiv\mathrm{id}$ (no update). In any closed-loop deployment,
actuation feedback shifts $\Pot^*$; since $\Pot$ cannot track, the
distributional mismatch grows at a rate lower-bounded by the rate of
environmental response to actuation.
\end{proposition}

\begin{proof}
The closed-loop case follows from
$\Pot^*(t+1)\leftarrow f(\Pot^*(t), \Act_t(\Pot(t)))$ with $\Pot$ fixed, so
the gap $W_2(\Pot,\Pot^*)$ grows at least as fast as the variation of $f$ in
its second argument whenever $\Act_t(\Pot)\neq\Act^*_t$.
\end{proof}

\begin{remark}[Interpretation]
FM-3 is the missing-adaptation regime. In passive settings it appears as a
failure to track external drift; in closed-loop settings it can amplify the
system's own feedback and produce self-reinforcing misalignment.
\end{remark}

These failure modes motivate the score construction in the next section. Once
the three ways of failing are separated, the natural diagnostic question is how
to summarize them in a single quantity without losing the identity of the active
bottleneck.
